
\section{Gossip min-mining}\label{sec:questions}
As motivation for our amnesiac approach, we want to look at other question strategies.

Various characters allow \player{}s to ask the \st questions.
A straightforward approach is to narrow down demon candidates with our question.
A common approach is to bisect the player set \playerset with questions such as

\begin{quote}
    \itshape
    Is player $X$ in the half of players that have the demon at their center?
\end{quote}
or
\begin{quote}
    \itshape
    Is the demon above the median age?
\end{quote}

What other questions can you think of?

\begin{question}\label{q:gossip-search}
    Let's discuss the following question
    \begin{quote}
        \itshape
        We give each player a color.
        Starting with the demon and continuing with the demon's next neighbor
        left, right, left right,
        we color half the number of players (rounded down) orange.
        The remaining players are colored green.
        Is player X orange?
    \end{quote}
    or in a less formal, more confusing and really boring way.
    \begin{quote}
        \itshape
        Is player X in the half of players that have the demon at their center?
    \end{quote}
\end{question}

\begin{observation}
    Let's look at two coloring examples for even and odd cases to
    \begin{seating}{Coloring of odd player count}[fig:q:coloring-odd]
    \seat{11}{1}{good}
    \seat{11}{2}{bad}
    \seat{11}{3}{bad}
    \seat{11}{4}{bad}[$D$]
    \seat{11}{5}{bad}
    \seat{11}{6}{bad}
    \seat{11}{7}{good}[$K$]
    \seat{11}{8}{good}
    \seat{11}{9}{good}
    \seat{11}{10}{good}
    \seat{11}{11}{good}
\end{seating}
    \begin{seating}{Coloring of even player count}[fig:q:coloring-even]
    \seat{12}{1}{good}
    \seat{12}{2}{bad}
    \seat{12}{3}{bad}
    \seat{12}{4}{bad}[$D$]
    \seat{12}{5}{bad}
    \seat{12}{6}{bad}
    \seat{12}{7}{bad}[$K$]
    \seat{12}{8}{good}
    \seat{12}{9}{good}
    \seat{12}{10}{good}
    \seat{12}{11}{good}
    \seat{12}{12}{good}
\end{seating}

    First, we observe that we know that either K
    or the \player(s) seated directly across from K are green and therefore not the demon.

    For \emph{\playercount is odd}, we learn two players that are not the demon

    Let's study $N=11$
\end{observation}

\begin{claim}
    Q\ref{q:gossip-search} clears exactly one player after one question, performs
    as good as any binary search after two or more questions and annoys the \st{ }nicely
\end{claim}

\begin{proof}
    We leave the proof as a nice exercise for you, dear reader.
\end{proof}

\begin{remark}[Ternary questions]
    The rulebook defines questions as ternary over the answer set
    ${\text{yes},\text{ no},\text{ i don't know}}$.
    As some redditor\todo{} showed, we can apply advanced question design to turn an ``I don't know'' in a ``We know more''
\end{remark}